\chapter{Hallucinations}

\section*{Definition}
Perception of sensory experiences (visual, auditory, tactile, olfactory, or gustatory) in the absence of external stimuli, arising from neurological or psychiatric dysfunction.

\section*{What Happens in Your Body}
Hallucinations involve aberrant activation of sensory cortex without corresponding external input. Auditory hallucinations (most common in schizophrenia) arise from misattribution of inner speech to external sources via dysfunction in frontotemporal networks. Visual hallucinations suggest organic pathology (delirium, dementia, seizures, migraine).

\section*{Causes}
\begin{itemize}
  \item Psychiatric disorders (schizophrenia, schizoaffective disorder, bipolar disorder)
  \item Delirium (especially in elderly or hospitalized people)
  \item Dementia (Lewy body dementia, Alzheimer disease)
  \item Substance use (LSD, psilocybin, PCP, cannabis)
  \item Alcohol withdrawal (delirium tremens)
  \item Seizures (temporal lobe epilepsy)
  \item Migraine with aura (scintillating scotomas, fortification spectra)
  \item Sensory deprivation (Charles Bonnet syndrome in visually impaired)
  \item Brain tumors or strokes affecting sensory pathways
  \item Metabolic derangements (electrolyte imbalance, hepatic encephalopathy)
\end{itemize}

\section*{When to See a Doctor}
See your doctor if:
\begin{itemize}
  \item Symptoms are severe or getting worse over time
  \item New-onset hallucinations with confusion, agitation, head injury, or risk of harm to self or others (emergency evaluation)
  \item You have other concerning symptoms such as fever, weight loss, or bleeding
\end{itemize}

\section*{Related Symptoms}
\begin{itemize}
  \item Delusions
  \item Confusion
  \item Agitation
  \item Sleep disturbance
  \item Anxiety
\end{itemize}
\textbf{Important:} If this symptom persists, your doctor will check for serious underlying causes.

\section*{Self-Care / Home Management}
\begin{itemize}
  \item Rest and avoid activities that make it worse.
  \item Maintain adequate hydration and a balanced diet.
  \item Over-the-counter remedies may provide symptomatic relief (consult a pharmacist or doctor).
  \item Monitor symptoms and seek professional advice if they persist or worsen.
  \item Avoid self-medicating with prescription medications without medical guidance.
\end{itemize}

\section*{How Your Doctor Will Evaluate You}
\begin{itemize}
  \item A thorough history and physical examination are the first steps in evaluation.
  \item Laboratory tests (blood work, urinalysis) may be ordered based on clinical suspicion.
  \item Imaging studies (X-ray, ultrasound, CT, MRI) may be indicated when structural pathology is suspected.
  \item Specialist referral is arranged when the diagnosis remains uncertain or requires specific expertise.
\end{itemize}

\section*{Treatment Options}
\begin{itemize}
  \item Treatment targets the underlying cause identified through diagnostic evaluation.
  \item Supportive care and symptom management are provided while awaiting definitive therapy.
  \item Medications, lifestyle modifications, and physical therapies may be prescribed as appropriate.
  \item Follow-up ensures treatment effectiveness and allows timely adjustments to the care plan.
\end{itemize}

\clearpage